\documentclass[conference]{IEEEtran}
\IEEEoverridecommandlockouts
\usepackage{cite}
\usepackage{balance}
\usepackage{amsmath,amssymb,amsfonts}
\usepackage{graphicx}
\usepackage{textcomp}
\usepackage{xcolor}
\usepackage{booktabs}
\usepackage{algorithm}
\usepackage{algorithmic}
\usepackage{tikz}
\usetikzlibrary{positioning,arrows.meta,fit,calc}
\usepackage[hidelinks]{hyperref}
\input{nums.tex}
\newtheorem{remark}{Remark}

\begin{document}

\title{One-Hot-Preserving QAOA for Flood Shelter\\
Assignment: Pairwise Capacity Encoding and Angle Transfer}

%<AUTHORS>
\author{
\IEEEauthorblockN{Nandhakumar T}
\IEEEauthorblockA{\textit{Department of Information Technology}\\
\textit{Sona College of Technology}\\
Salem, Tamil Nadu, India\\
nandhakumartm93@gmail.com}
\and
\IEEEauthorblockN{Balakarthikeyan B S}
\IEEEauthorblockA{\textit{Department of Artificial Intelligence and Data Science}\\
\textit{Sona College of Technology}\\
Salem, Tamil Nadu, India\\
balasan2626@gmail.com}
\and
\IEEEauthorblockN{Gowthaman J}
\IEEEauthorblockA{\textit{Department of Artificial Intelligence and Data Science}\\
\textit{Sona College of Technology}\\
Salem, Tamil Nadu, India\\
gowthamanjanadhran@gmail.com}
\and
\IEEEauthorblockN{Abhishek S}
\IEEEauthorblockA{\textit{Department of Artificial Intelligence and Data Science}\\
\textit{Sona College of Technology}\\
Salem, Tamil Nadu, India\\
abisheksengodan18@gmail.com}
}
%</AUTHORS>

\maketitle

\begin{abstract}
Flood shelter assignment is a small capacity-constrained assignment problem that is often proposed for quantum optimization. We study the Quantum Approximate Optimization Algorithm (QAOA) by exact classical simulation on instances from a flood digital twin and from the road networks of \CityNCities{} Indian cities. Capacity enters as approximate pairwise penalties, so a classical verifier checks every sampled plan. At 12 qubits and depth 4, a ring-XY mixer that keeps each group assigned to one shelter placed a mean probability of \SeedXYMean{} on the optimal plan, against at most \SeedXMax{} for the transverse-field mixer at the same CNOT count. One angle set per depth, chosen on development instances with their optima, reached \SelPopt{} on held-out twin seeds against \SelTrained{} for per-instance training; chosen by energy or CVaR, it fell below per-instance training in leave-one-out. A rule a device can apply, the expected cost of the best verified plan among \Shots{} shots, picked the same set as the optima in all \LcFolds{} leave-one-city-out folds over four disruption patterns. On the held-out cities it left the optimal-plan probability unchanged (\LcVOpt{} against \LcTrained) but raised the chance that \Shots{} verified samples contain the optimum from \LcTrainedSucc\% to \LcVSucc\%. Up to \ScQubitsMax{} qubits the optimal-plan probability fell about \ScDecay-fold per added group as capacity pushed probability onto infeasible plans, and an exact capacity penalty raised feasibility but not optimality. Exhaustive search solves the 12-qubit instances in about a millisecond; we claim no quantum advantage.
\end{abstract}

\begin{IEEEkeywords}
QAOA, XY mixer, parameter transfer, QUBO, evacuation planning
\end{IEEEkeywords}

\section{Introduction}
When a flood closes roads and fills shelters, an emergency officer must decide again which shelter each evacuation group should go to, trading off travel time, flood exposure and the seats left. As a binary program this is a capacity-constrained assignment problem, the kind that is routinely encoded as a quadratic unconstrained binary optimization (QUBO) problem \cite{Lucas2014,Glover2019} and offered as an application for near-term quantum optimization \cite{Cadelina2024,Abbas2024}.

Prior proposals rarely test whether the circuit keeps the constraints or only penalizes violations, whether the angles must be retrained after every disruption, and how the results compare with classical methods and with a sampler followed by a feasibility check. The original QAOA \cite{Farhi2014} explores all $2^n$ bitstrings, most of which are not plans, while the quantum alternating operator ansatz \cite{Hadfield2019} keeps the state in a feasible subspace \cite{WangZ2020,Fuchs2022,Bartschi2020}; mixer studies, however, treat equality constraints, and capacity is an inequality.

We address them on instances generated by the digital twin of a flood decision-support prototype. Our contributions are:
\begin{enumerate}
\item A shelter-assignment QUBO in which capacity is encoded by pairwise products instead of slack qubits, with a cost shift and scaling that place all instances on one energy range; we show that it admits infeasible zero-penalty plans and compare it with an exact penalty over minimal overloading sets, integer programming, exhaustive search, a greedy rule and simulated annealing.
\item A comparison of transverse-field, ring-XY, Grover and feasible-subspace mixers under mean-energy and conditional value-at-risk (CVaR) training, with \NSeeds{} optimizer seeds for the X and XY mixers, with Holm-corrected tests and per-layer CNOT counts.
\item An angle-transfer study with selection rules with and without optima, including a verified-cost rule a device can apply, tested on held-out twin seeds, on city networks under four disruption patterns with leave-one-city-out selection, and up to \ScQubitsMax{} qubits, with an ablation in which transfer failed without the shared energy scale.
\end{enumerate}

\section{Related Work}
\textit{Formulations.} Lucas gave Ising forms for many NP-hard problems \cite{Lucas2014}, and Glover \textit{et al.} describe how assignment and capacity constraints become quadratic penalties, with slack variables for inequalities \cite{Glover2019}. Evacuation shelter allocation has been solved by integer programming \cite{Sherali1991}, and disaster resource allocation by quantum annealing \cite{Cadelina2024}.

\textit{Mixers.} QAOA alternates a cost unitary and a mixing unitary for $p$ layers \cite{Farhi2014,Blekos2024}. Hadfield \textit{et al.} generalized the mixer to preserve hard constraints \cite{Hadfield2019}. Wang \textit{et al.} analyzed XY mixers for one-hot constraints \cite{WangZ2020}, Fuchs \textit{et al.} built constraint-preserving \cite{Fuchs2022} and subspace mixers \cite{Fuchs2024}, and Grover mixers move the complexity into state preparation \cite{Bartschi2020}. XY mixers have also served routing \cite{doCarmo2025} and rostering \cite{Gupta2026}.

\textit{Training and transfer.} Angles can be interpolated across depths \cite{Zhou2020}, CVaR objectives favor good samples \cite{Barkoutsos2020}, and warm starts use classical solutions \cite{Egger2021}. The objective value at fixed angles concentrates across typical instances \cite{Brandao2018}, optimal angles transfer between random graphs \cite{Galda2021}, and for weighted MaxCut a rescaling of the weights lets one parameter set serve many instances \cite{Shaydulin2023}. On hardware, QAOA performance falls with size and depth \cite{Harrigan2021}, and practical advantage has not been shown \cite{Abbas2024}.

\textit{Positioning.} The ring-XY mixer is that of \cite{WangZ2020}, and the transfer studies \cite{Brandao2018,Galda2021,Shaydulin2023} use unconstrained MaxCut with the standard mixer. We add a capacity inequality, as a pairwise or exact penalty, next to a one-hot-preserving mixer, compare selection rules with and without optima, and test on held-out instances and real road networks with a verifier in the loop.

\section{Problem Formulation}
\subsection{Instances from the Digital Twin}
The twin models a corridor with three population zones, nine junctions and three shelters joined by 22 roads. Every road carries a travel time, a flood-depth index and a congestion level drawn from a seed. Shelters have capacities of 120, 90 and 70 people with an initial occupancy of 15--45\%, and groups have 8--22 people. Each instance applies the disruption used in the twin's routing benchmark: road R5 closes and Shelter A receives arrivals equal to half its capacity. For group $g$ and shelter $s$ we take the best open path under the twin's flood- and congestion-weighted edge cost and define
\begin{equation}
c_{gs}=T_{gs}+5H_{gs},
\end{equation}
where $T_{gs}$ is travel time in minutes and $H_{gs}$ the mean flood index along the path. The pair $(g,s)$ is \emph{forbidden} if no open path exists, the shelter has failed, or the group alone exceeds the remaining seats $R_s$. With $G$ groups, binary variables $x_{gs}$ give $n=3G$ qubits. We write $m_g$ for the number of people in group $g$. A plan is feasible if every group has exactly one shelter, no forbidden pair is used and
\begin{equation}
\textstyle\sum_g m_g x_{gs}\le R_s\quad\text{for each shelter } s.
\end{equation}

\subsection{QUBO with Pairwise Capacity}
An inequality normally needs slack qubits \cite{Glover2019}. We use a pairwise encoding instead. Let $\mathcal{O}=\{(g,h,s): m_g+m_h>R_s\}$ be the pairs of groups that cannot share shelter $s$ and $\mathcal{F}$ the forbidden pairs. With shifted costs $c'_{gs}=c_{gs}-\min_{s'}c_{gs'}$ and penalty weight $\mu$,
\begin{multline}
E(\mathbf{x})=\sum_{g,s}c'_{gs}x_{gs}+\mu\Big[\sum_g\Big(1-\sum_s x_{gs}\Big)^2\\
+\sum_{(g,s)\in\mathcal{F}}x_{gs}+\sum_{(g,h,s)\in\mathcal{O}}x_{gs}x_{hs}\Big].
\label{eq:qubo}
\end{multline}

\begin{remark}
With $\mu=1+\sum_g\max_s c'_{gs}$, every bitstring that breaks a one-hot, forbidden-pair or pairwise-capacity term has higher energy than every bitstring that breaks none of them, since $c'_{gs}\ge0$ bounds the cost of a bitstring without violations by $\mu-1$ and each violation adds at least $\mu$. The bound no longer holds for the reduced weights $\lambda\mu$ with $\lambda<1$ used in the penalty sweep of Section~\ref{sec:mixres}.
\end{remark}

On one-hot states the shift subtracts the constant $\sum_g\min_s c_{gs}$, so the optimal plan is unchanged. We divide the Hamiltonian by $\Delta=\mu-1$, which maps the energies of all penalty-free plans into $[0,1]$ for every instance. The pairwise encoding misses overloads by three or more groups whose pairs all fit, so a penalty-free plan can still be infeasible, and every sampled plan is re-checked exactly before use.

An exact penalty also needs no slack qubits. Let $\mathcal{M}$ hold the minimal overloading sets $(T,s)$: groups $T$ allowed at $s$ with $\sum_{g\in T}m_g>R_s$, while every proper subset fits. Replacing the last sum in (\ref{eq:qubo}) by
\begin{equation}
\textstyle\sum_{(T,s)\in\mathcal{M}}\prod_{g\in T}x_{gs}
\label{eq:hubo}
\end{equation}
leaves no capacity penalty exactly when no shelter is overloaded, because an overloaded shelter always holds a minimal set; the pairs of $\mathcal{O}$ are the sets of size two. A product of $k$ variables expands into $2^k-1$ Pauli-$Z$ strings, and a string of weight $w$ costs $2(w-1)$ CNOTs.

\section{QAOA Variants and Training}
For depth $p$ and angles $\boldsymbol\theta=(\boldsymbol\gamma,\boldsymbol\beta)$, QAOA prepares
\begin{equation}
|\psi(\boldsymbol\theta)\rangle=\prod_{k=1}^{p}e^{-i\beta_k B}\,e^{-i\gamma_k \hat E/\Delta}\,|\psi_0\rangle,
\end{equation}
where $\hat E$ is diagonal with entries $E(\mathbf{x})$. Substituting $x_q=(1-Z_q)/2$ turns (\ref{eq:qubo}) into an Ising Hamiltonian with one $R_Z$ rotation per qubit and one $ZZ$ rotation, two CNOTs, per non-zero coupling.

\textit{Transverse-field (X) mixer.} $B=\sum_q X_q$ and $|\psi_0\rangle=|+\rangle^{\otimes n}$; constraints act only through the penalties. A cost layer has 12 one-hot couplings and on average \OvPairs{} capacity couplings, a mean of \CnotX{} CNOTs with all-to-all connectivity; the mixer needs none.

\textit{Ring-XY mixer.} For each group $g$ with qubits $a_g,b_g,c_g$,
\begin{equation}
B_{\mathrm{XY}}=\sum_{g}\ \sum_{(u,v)\in\{(a_g,b_g),(b_g,c_g),(c_g,a_g)\}}\tfrac12\big(X_uX_v+Y_uY_v\big).
\label{eq:xy}
\end{equation}
Each term swaps amplitude between $|01\rangle$ and $|10\rangle$, so a one-hot group stays one-hot \cite{WangZ2020}. The start state $|W_3\rangle^{\otimes G}$, with $|W_3\rangle=(|100\rangle+|010\rangle+|001\rangle)/\sqrt3$, is the uniform superposition of the $3^G$ one-hot plans \cite{Bartschi2019}. The one-hot penalty is then constant and dropped; the three ring terms are applied in a fixed order. Each XY rotation compiles to two CNOTs \cite{Vatan2004}, so a layer has \CnotXY{} CNOTs, the same count as the X variant. We prepare $|W_3\rangle$ with an $R_Y$ rotation, a controlled Hadamard (one CNOT) and two CNOTs, which we checked numerically; this adds \CnotW{} CNOTs once.

\textit{Grover mixer.} $B_{\mathrm{GM}}=|\psi_0\rangle\langle\psi_0|$ with the same start state, so $e^{-i\beta B_{\mathrm{GM}}}=I+(e^{-i\beta}-1)|\psi_0\rangle\langle\psi_0|$ \cite{Bartschi2020}. A layer needs the inverse preparation, a phase on $|0\rangle^{\otimes n}$ controlled by all $n$ qubits, and the preparation. With the textbook ancilla construction \cite{Nielsen2010} the controlled phase takes $2(n-2)$ six-CNOT Toffolis, one controlled phase and $n-2=\NAnc$ ancillas (\CnotMCP{} CNOTs), so a Grover layer has \CnotGM{} CNOTs.

\textit{Feasible-subspace mixer.} As a reference bound, not a usable circuit, a Grover mixer starts from the uniform superposition of the feasible plans, which requires the feasible set \cite{Bartschi2020,Fuchs2024}.

\textit{Objectives and training.} We minimize either $\langle\hat E\rangle/\Delta$ or the CVaR with $\alpha=0.1$ \cite{Barkoutsos2020}, the only level we tried; it was not tuned. Sorting bitstrings by energy, $E_{(1)}\le E_{(2)}\le\dots$, with probabilities $\pi_{(k)}$ and cumulative mass $\Pi_k=\sum_{j\le k}\pi_{(j)}$,
\begin{equation}
\mathrm{CVaR}_\alpha=\frac1\alpha\sum_k E_{(k)}\,\min\!\big(\pi_{(k)},\ \max(0,\ \alpha-\Pi_{k-1})\big).
\label{eq:cvar}
\end{equation}
Every variant is trained from scratch with the same procedure: COBYLA \cite{Powell1994} with at most 250 evaluations per start, depth by depth. At each depth $p$ we start from two random points ($\gamma_k\sim U(0,\pi)$, $\beta_k\sim U(0,\pi/2)$) and, for $p>1$, from the optimum of depth $p-1$ interpolated to $p$ layers \cite{Zhou2020}, and keep the start with the lowest objective; training never uses the optimum.

\subsection{Angle Transfer}\label{sec:transfer}
Since the scaled Hamiltonians of all instances share the energy range $[0,1]$, we test whether trained angles carry over. Let $\Theta_p=\{\boldsymbol\theta_j^{(p)}:j\in\mathcal{D}\}$ be the angles trained on the development instances $\mathcal{D}$. We fix one set $\boldsymbol\theta^{\star}_p\in\Theta_p$ by one of three rules:
\begin{equation}
\boldsymbol\theta^{\star}_p=\arg\max_{\boldsymbol\theta\in\Theta_p}\ \tfrac{1}{|\mathcal{D}|}\textstyle\sum_{j\in\mathcal{D}}P_{\mathrm{opt}}(j;\boldsymbol\theta),
\label{eq:transfer}
\end{equation}
or the set with the lowest mean of $\langle\hat E\rangle/\Delta$, or with the lowest mean CVaR, over $\mathcal{D}$. Rule (\ref{eq:transfer}) needs the exact optimum of every development instance, from exhaustive search; it is an oracle on the development set, though it uses nothing from the instances it is applied to. The other two rules use only quantities a device can estimate. We choose among whole trained sets because their coordinate-wise median, which mixes basins of the periodic angle space, gave $P_{\mathrm{opt}}=\LooMedianFour$ in leave-one-out.

\section{Experimental Setup}\label{sec:setup}
\textit{Instances.} The development set has \NInst{} instances from twin seeds $42+17i$, $i=0,\dots,29$, with the first four groups ($G=4$, $n=12$). All are feasible, with on average \MeanFeasAssign{} feasible plans among the 81 one-hot plans, and in \NBinding{} of them capacity changes the optimal plan (\emph{capacity-binding}). The held-out set has \HoN{} further seeds, $i=30,\dots,59$, of the same disruption pattern, used only for the final transfer test. We also built \CityAll{} instances (\CityN{} feasible, \CityBinding{} capacity-binding) on the OpenStreetMap road graphs of \CityNCities{} flood-prone Indian cities, with each city's mapped zones, shelters and road flood-exposure proxy. Group sizes and seats follow the twin, as demand is not mapped; the disruption closes the most flood-exposed road on a zone--shelter path and fills half of the first shelter's capacity.

\textit{Metrics.} Simulations store all $2^n$ amplitudes in double precision, with fixed seeds; code and stored results will be released after review. $P_{\mathrm{opt}}$ is the probability of measuring an optimal plan, a feasible plan of lowest true cost (not always the minimizer of (\ref{eq:qubo})), and $P_{\mathrm{feas}}$ that of a feasible plan. A sampler draws $N$ plans and returns the cheapest one that passes an exact feasibility check. It returns the optimum whenever the optimum is among the samples, which has probability success@$N=1-(1-P_{\mathrm{opt}})^N$, assuming independent shots. We also report the expected relative gap $(c-c^\star)/c^\star$ of the returned plan, given that a feasible plan was drawn, for $N\in\{8,32,128\}$, computed exactly from the output distribution. An evaluation is one estimate of the training objective (on hardware, a batch of circuit runs).

\textit{Statistics.} Paired comparisons use the two-sided Wilcoxon signed-rank test on per-instance values, also for success@$N$. We apply the Holm correction within each family of tests: the four depths of one mixer comparison, the \SelFamily{} rows of Table~\ref{tab:select}, the four city comparisons, the \DsFamily{} tests of Table~\ref{tab:dis}, the four leave-one-city-out tests, the \ScFamily{} size comparisons and the four tests of Table~\ref{tab:exact}. Effect sizes are median paired differences. Intervals for variants run with \NSeeds{} optimizer seeds use a two-level bootstrap that resamples instances and then one seed per instance. Because the city instances cluster by city, we also apply a sign test to the per-city mean differences.

\textit{Baselines.} Uniform sampling over the 81 one-hot plans ($P_{\mathrm{opt}}=\UniOneHot$) and over the feasible plans ($P_{\mathrm{opt}}=\UniFeas$); exhaustive search; the integer program solved by HiGHS \cite{Huangfu2018}; a greedy rule that places the largest group first in its cheapest shelter with seats left; and simulated annealing (SA) \cite{Kirkpatrick1983} with 20 restarts of 200 sweeps, either with bit flips on (\ref{eq:qubo}) or with moves that reassign one group, which keep the plan one-hot. For the X mixer we swept the penalty $\lambda\mu$ with $\lambda\in\{0.25,0.5,1,2\}$.

\section{Results}
\subsection{Mixers, Objectives and Penalties}\label{sec:mixres}
Table~\ref{tab:mixers} and Fig.~\ref{fig:popt} compare the mixers. Over \NSeeds{} seeds the XY mixer reached $P_{\mathrm{opt}}=\SeedXYMean$ at $p=4$ (standard deviation over seeds \SeedXYSd), against \SeedXMin--\SeedXMax{} for the X mixer, and it was better on \StXYvXWinsMin--\StXYvXWinsMax{} of \NPaired{} instances at every depth (Holm $p\le\StXYvXMaxP$) and for every seed. It kept $P_{\mathrm{feas}}$ between \XYPfeasMin{} and \XYPfeasMax, while the X mixer started at \QXMeanOnePfeas. As every sample is verified, the fairer figure is the chance that a feasible sample is optimal, $P_{\mathrm{opt}}/P_{\mathrm{feas}}$: \CondXY{} for XY against \CondX{} for X.

\begin{table}[t]
\caption{Mixers and objectives at depth 4 (\NPaired{} development instances, 12 qubits)}
\label{tab:mixers}
\centering
\footnotesize
\setlength{\tabcolsep}{2.6pt}
\begin{tabular}{lcccccc}
\toprule
Variant & $P_{\mathrm{opt}}$ & 95\% CI & $P_{\mathrm{feas}}$ & s@\Shots\,\% & Evals & CNOT\\
\midrule
\input{tab_mixers.tex}
\bottomrule
\multicolumn{7}{l}{\scriptsize s@\Shots: success@\Shots. X, XY: mean over \NSeeds{} seeds; CI: two-level bootstrap.}\\
\multicolumn{7}{l}{\scriptsize Evals: per instance, depths 1--4. CNOT: per layer, all-to-all.}\\
\end{tabular}
\end{table}

\begin{figure}[t]
\centering
\includegraphics[width=0.92\columnwidth]{fig_popt.pdf}
\caption{Probability of the optimal plan (left, different $y$-range) and of any feasible plan (right) against depth, optimizer seed 0. Dotted: uniform one-hot.}
\label{fig:popt}
\end{figure}

Trained once (seed 0), the Grover mixer reached \GMMeanFour{} against \QXYMeanFourPopt{} for XY with the same seed. The difference was not significant after Holm correction at any depth ($p\ge\StXYvGMMinP$), although XY led on \GmWinsFour{} of \NPaired{} instances at $p=4$ (mean difference \GmDiffFour, bootstrap interval \GmDiffLo--\GmDiffHi, unadjusted $p=\GmRawP$). A Grover layer costs \CnotGM{} CNOTs and \NAnc{} ancillas against \CnotXY. CVaR training lowered XY's mean $P_{\mathrm{opt}}$ but raised success@\Shots{} from \QXYMeanFourSucc\% to \QXYCvarFourSucc\% (seed 0), lifting the worst instance from \XYMin{} to \XYCvarMin. Lowering the X mixer's penalty to $\lambda=\PenBestLam$ raised its $P_{\mathrm{opt}}$ \PenGain-fold to \PenBestPopt, at $P_{\mathrm{feas}}=\PenBestFeas$. A mixer restricted to feasible plans, a bound on what the capacity penalty costs, reached \FgFour, \FgOverXYZero{} times XY with the same seed ($p=\FgPFour$).

\subsection{Angle Transfer to Held-Out Seeds}
\begin{table*}[t]
\caption{Angles at $p=4$ on the \HoN{} held-out seeds: per-instance training against one fixed angle set chosen on the development set}
\label{tab:select}
\centering
\footnotesize
\setlength{\tabcolsep}{3.3pt}
\begin{tabular}{llccccccc}
\toprule
Variant & Angles & Needs optima & Held-out $P_{\mathrm{opt}}$ & Median paired diff. & Wins of \HoN & Holm $p$ & Leave-one-out $P_{\mathrm{opt}}$ (dev.) & s@\Shots\,\%\\
\midrule
\input{tab_select.tex}
\bottomrule
\multicolumn{9}{l}{\scriptsize Fixed sets vs.\ per-instance training of the same variant (Wilcoxon signed-rank, Holm over \SelFamily{} rows). Per-instance rows: development-set mean in the leave-one-out column.}\\
\multicolumn{9}{l}{\scriptsize For XY, CVaR, both rules chose the same set.}\\
\end{tabular}
\end{table*}

Table~\ref{tab:select} gives the held-out test. The set of rule (\ref{eq:transfer}) reached $P_{\mathrm{opt}}=\SelPopt$ against \SelTrained{} for per-instance training with \HoCumFour{} evaluations (better on \SelPoptWins{} of \HoN, Holm $p=\SelPoptP$). The energy and CVaR sets also won on these seeds (\SelEn{} and \SelCv), but in leave-one-out on the development set they gave \SelEnLoo{} and \SelCvLoo, below the \SelDevTrained{} of per-instance training, while rule (\ref{eq:transfer}) gave \SelPoptLoo; CVaR-trained angles behaved alike. The held-out seeds had more capacity-binding instances (\HoNBind{} against \NBinding), which lowered per-instance training there; the gain came mostly from them (\HoBindTrained{} to \HoBindFixed, against \HoNonBindTrained{} to \HoNonBindFixed{} on the others). For the X mixer the fixed set gained little (\SelXPopt).

Per-instance training reached a lower mean energy than the fixed set (\EnTrained{} against \EnFixed), but across the \CandN{} trained sets mean energy and mean $P_{\mathrm{opt}}$ had a rank correlation of only \RhoEnPopt. Stronger optimizers stayed below the fixed set (best of five COBYLA seeds \SelBestFive, L-BFGS-B \SelLbfgs), while COBYLA started from the fixed angles reached \StrWarm{} ($p=\StrWarmP$), the best result here. A set's development $P_{\mathrm{opt}}$ predicted its held-out $P_{\mathrm{opt}}$ (Spearman \RhoDevHeld), though only \CandBeat{} of the \CandN{} sets beat per-instance training. Without the energy scaling the fixed angles reached only \NormNoneFixed; scaling by the standard deviation of penalty-free energies gave \NormStdFixed, so a common scale is needed but ours is not special. The one-time cost, \OffTotal{} evaluations, equals per-instance training on about \OffBreakEven{} disruptions. With 8 verified samples the fixed angles came within \GapFxEight\% of the optimal cost on average, per-instance training within \GapTrEight\%.

\subsection{Transfer to Real Road Networks}
With angles fixed on the twin, over the \CityN{} city instances rule (\ref{eq:transfer}) gave $P_{\mathrm{opt}}=\CityFixed$ against \CityTrained{} (Holm $p=\CityPHolm$), leading in only \CityPos{} of \CityK{} cities, while success@\Shots{} rose from \CitySuccTrained\% to \CitySuccFixed\% ($p=\CitySuccPHolm$), in all \CitySuccPos{} cities (sign test $p=\CitySuccSignP$).

\textit{Other disruptions.} We rebuilt the city instances under three disruptions the twin never saw: two roads closed; one road closed and the first shelter failed; one road closed and 70\% of the second shelter's seats cut (\DsNTotal{} feasible instances, Table~\ref{tab:dis}). No fixed set raised $P_{\mathrm{opt}}$ significantly, and two lost it: rule (\ref{eq:transfer}) when a shelter failed (\DsShelterO{} against \DsShelterTrained, Holm $p=\DsShelterOP$) and the energy rule when two roads closed ($p=\DsRoadsEP$). Pooled, success@\Shots{} still rose from \DsAllTrainedSucc\% to \DsAllOSucc\% for rule (\ref{eq:transfer}) ($p=\DsAllOSuccP$) and \DsAllESucc\% for the energy rule ($p=\DsAllESuccP$).

\begin{table}[t]
\caption{City instances under other disruptions (XY mixer, $p=4$): per-instance training against the fixed sets of rule (\ref{eq:transfer}) (fix-O) and the energy rule (fix-E)}
\label{tab:dis}
\centering
\footnotesize
\setlength{\tabcolsep}{1.6pt}
\begin{tabular}{lccccccc}
\toprule
 & \multicolumn{5}{c}{$P_{\mathrm{opt}}$} & \multicolumn{2}{c}{s@\Shots\,\%}\\
\cmidrule(lr){2-6}\cmidrule(lr){7-8}
Pattern & trained & fix-O & Holm $p$ & fix-E & Holm $p$ & trained & fix-O\\
\midrule
\input{tab_dis.tex}
\bottomrule
\end{tabular}
\end{table}

\textit{Choosing a set without optima.} A device can estimate, for each candidate set, the expected cost of the cheapest verified plan among \Shots{} shots, with the greedy plan when no shot is valid, relative to the greedy plan's cost. Over the \LcCand{} development sets this score ranked the sets like their mean $P_{\mathrm{opt}}$ (Spearman \LcRhoV, against \LcRhoE{} for mean energy). We chose a set with it on five cities and tested it on the sixth, over all four disruption patterns (\LcN{} instances); in all \LcFolds{} folds it picked the set that the optima would have picked. On the held-out cities $P_{\mathrm{opt}}$ was \LcVOpt{} against \LcTrained{} for per-instance training (Holm $p=\LcVOptP$), and success@\Shots{} rose from \LcTrainedSucc\% to \LcVSucc\% ($p=\LcVSuccP$), in all \LcVSuccCities{} cities; the energy rule reached \LcESucc\%.

\subsection{Larger Instances and Gate Noise}
The ring-XY circuit stays in the $3^G$ one-hot amplitudes, so we simulated it there; on \SubCheckRuns{} random 12- and 15-qubit circuits this agreed with the full statevector to $\SubCheckErr$. On \ScN{} new twin instances per size from \ScQubitsMin{} to \ScQubitsMax{} qubits (Table~\ref{tab:size}), the set of rule (\ref{eq:transfer}) beat per-instance training at every size (Holm $p\le\ScOMaxP$), but that training itself fell below $P_{\mathrm{opt}}=0.01$ on at least \ScBigTrainedLowMin{} of \ScN{} instances from \ScBigQ{} qubits on; the energy-rule set was significant only at \ScESigQ{} qubits. The fixed set's $P_{\mathrm{opt}}$ fell about \ScDecay-fold per added group, to \ScNineO{} at \ScQubitsMax{} qubits, still a median \ScNineOGain{} times uniform one-hot sampling. Capacity explains most of the drop: from \ScBindAllFromQ{} qubits on every instance was capacity-binding and at most \ScFeasMaxBig{} of the probability lay on feasible plans.

\begin{table}[t]
\caption{Fixed 12-qubit angles on \ScN{} twin instances per size $n$ (qubits), XY mixer, $p=4$; fix-O: set of rule (\ref{eq:transfer}), fix-E: energy rule; $P_{\mathrm{feas}}$ of fix-O}
\label{tab:size}
\centering
\footnotesize
\setlength{\tabcolsep}{2.2pt}
\begin{tabular}{ccccccccc}
\toprule
 & \multicolumn{5}{c}{$P_{\mathrm{opt}}$} & & \multicolumn{2}{c}{s@\Shots\,\%}\\
\cmidrule(lr){2-6}\cmidrule(lr){8-9}
$n$ & trained & fix-O & Holm $p$ & fix-E & Holm $p$ & $P_{\mathrm{feas}}$ & trained & fix-O\\
\midrule
\input{tab_beyond.tex}
\bottomrule
\end{tabular}
\end{table}

In a trajectory model where each two-qubit gate fails with probability $\epsilon=0.01$, a CNOT error near 0.5\% and close to the \SycIso--\SycSim\% of a superconducting processor \cite{Arute2019}, XY trained per instance kept $P_{\mathrm{opt}}=\TrajXYOne$ on \TrajInst{} development instances, \TrajRatioOne{} times the X mixer's \TrajXOne.

\subsection{Encoding Exactness and Classical Baselines}
In all \ZeroPenInfeasAll{} development instances some infeasible one-hot plan had zero penalty, yet the minimizer of (\ref{eq:qubo}) was the optimum in \QuboExact\% of them. Exhaustive search took \ExactMs~ms, the greedy rule found the optimum in \GreedyOpt\%, and both annealers visited the optimum every time, so with a verifier they found it in \SAohVer\%. On the \ExN{} city instances the minimal-set penalty (\ref{eq:hubo}) was exact, checked over all 81 plans, while the pairwise penalty admitted \ExMissPlans{} infeasible penalty-free plans in \ExMissInst{} instances; exhaustive search took \ExExactMs~ms and HiGHS \ExMilpMs~ms, both exact. With per-instance training (Table~\ref{tab:exact}, an independent rerun), the exact penalty raised $P_{\mathrm{feas}}$ from \ExTrFeasPair{} to \ExTrFeasHubo{} (Holm $p=\ExTrFeasP$, all \ExFeasCities{} cities) but not $P_{\mathrm{opt}}$ (\ExTrOptPair{} against \ExTrOptHubo, $p=\ExTrOptP$), and the fixed angles lost $P_{\mathrm{opt}}$ on it (\ExFxOptPair{} to \ExFxOptHubo, $p=\ExFxOptP$). Its higher-order terms cost \ExCnotHubo{} CNOTs per layer (up to \ExCnotHuboMax) against \ExCnotPair, and a slack encoding would add \ExSlackMin{} to \ExSlackMax{} qubits.

\begin{table}[t]
\caption{Pairwise vs.\ exact capacity penalty, \ExN{} city instances, $p=4$. miss: instances with an infeasible penalty-free plan; min.: minimizer optimal}
\label{tab:exact}
\centering
\footnotesize
\setlength{\tabcolsep}{2.2pt}
\begin{tabular}{lccccccc}
\toprule
 & & & & \multicolumn{2}{c}{trained} & \multicolumn{2}{c}{fix-O}\\
\cmidrule(lr){5-6}\cmidrule(lr){7-8}
Penalty & miss & min.\,\% & CNOT & $P_{\mathrm{feas}}$ & $P_{\mathrm{opt}}$ & $P_{\mathrm{feas}}$ & $P_{\mathrm{opt}}$\\
\midrule
\input{tab_exact.tex}
\bottomrule
\end{tabular}
\end{table}

\section{Discussion and Limitations}
For hardware experiments the results suggest the following. First, the one-hot constraint belongs in the mixer: at the same CNOT count, ring-XY beat the transverse-field mixer on almost every instance and seed. Capacity is not solved by any variant we tried. The feasible-subspace mixer raised $P_{\mathrm{opt}}$ \FgOverXYZero-fold but needs the feasible set; once capacity bound on every instance, ring-XY put almost all its probability on infeasible plans; and the exact penalty raised feasibility at 12 qubits but not $P_{\mathrm{opt}}$, which fell for the fixed angles, for more than twice the CNOTs. We plan to test a mixer that also preserves capacity. Second, transfer needs a common energy scale and a good selection rule. With development-set optima one fixed set beat per-instance training on held-out twin seeds, for a one-time cost of \OffTotal{} evaluations; the verified-cost rule made the same choice without optima on unseen cities. Under new disruptions no fixed set raised $P_{\mathrm{opt}}$, so the gain lies in the reliability of verified sampling, not in the optimal-plan probability. Third, the pairwise penalty admits infeasible penalty-free plans, so every sample must be verified.

The exact penalty was tested on the 12-qubit city instances only. All twin instances share one disruption pattern; the other disruptions were tested on the 12-qubit city instances only, and the city instances use sampled demand. The noise model charges errors per logical gate, ignores single-qubit and correlated errors and was run at 12 qubits; no circuit was compiled for hardware. $P_{\mathrm{opt}}$ depends on the training budget, and the Grover test cannot prove equivalence.

\section{Conclusion}
Ring-XY gave QAOA a significantly higher $P_{\mathrm{opt}}$ than the transverse-field mixer at equal CNOT count. One fixed angle set per depth, chosen with development-set optima or, without them, by the expected cost of the best verified plan, raised the chance that \Shots{} verified samples contain the optimum on unseen cities and disruptions, although it did not raise $P_{\mathrm{opt}}$ there; rules based on energy were weaker. Capacity is the open problem: the pairwise penalty needs a verifier, and an exact penalty raised feasibility but not $P_{\mathrm{opt}}$. At these sizes exhaustive search is faster, so the results give no evidence of advantage.

\section*{Acknowledgment}
We thank the OpenStreetMap contributors and Open-Meteo for open data used by the twin.
In line with the IEEE policy on AI-generated content, the authors disclose that Claude (Anthropic) was used to help edit and restructure the text of all sections and to draft the simulation, experiment and plotting scripts. All experiments are reproducible with fixed seeds. The authors checked every number, claim and reference and take full responsibility for the content.

\balance
\begin{thebibliography}{00}
\bibitem{Lucas2014} A. Lucas, ``Ising formulations of many NP problems,'' \textit{Front. Phys.}, vol. 2, Art. no. 5, 2014, doi: 10.3389/fphy.2014.00005.
\bibitem{Glover2019} F. Glover, G. Kochenberger, and Y. Du, ``Quantum bridge analytics I: A tutorial on formulating and using QUBO models,'' \textit{4OR}, vol. 17, no. 4, pp. 335--371, 2019, doi: 10.1007/s10288-019-00424-y.
\bibitem{Cadelina2024} A. R. Cadeli\~{n}a \textit{et al.}, ``D-Wave implementation of quantum annealing for optimal resource allocation in disaster response operation of Marikina City,'' \textit{ECTI Trans. Comput. Inf. Technol.}, vol. 18, no. 1, pp. 24--33, 2024, doi: 10.37936/ecti-cit.2024181.253577.
\bibitem{Abbas2024} A. Abbas \textit{et al.}, ``Challenges and opportunities in quantum optimization,'' \textit{Nat. Rev. Phys.}, vol. 6, pp. 718--735, 2024, doi: 10.1038/s42254-024-00770-9.
\bibitem{Farhi2014} E. Farhi, J. Goldstone, and S. Gutmann, ``A quantum approximate optimization algorithm,'' arXiv:1411.4028, 2014.
\bibitem{Hadfield2019} S. Hadfield, Z. Wang, B. O'Gorman, E. G. Rieffel, D. Venturelli, and R. Biswas, ``From the quantum approximate optimization algorithm to a quantum alternating operator ansatz,'' \textit{Algorithms}, vol. 12, no. 2, Art. no. 34, 2019, doi: 10.3390/a12020034.
\bibitem{WangZ2020} Z. Wang, N. C. Rubin, J. M. Dominy, and E. G. Rieffel, ``XY mixers: Analytical and numerical results for the quantum alternating operator ansatz,'' \textit{Phys. Rev. A}, vol. 101, no. 1, Art. no. 012320, 2020, doi: 10.1103/PhysRevA.101.012320.
\bibitem{Fuchs2022} F. G. Fuchs, K. O. Lye, H. M. Nilsen, A. J. Stasik, and G. Sartor, ``Constraint preserving mixers for the quantum approximate optimization algorithm,'' \textit{Algorithms}, vol. 15, no. 6, Art. no. 202, 2022, doi: 10.3390/a15060202.
\bibitem{Bartschi2020} A. B\"{a}rtschi and S. Eidenbenz, ``Grover mixers for QAOA: Shifting complexity from mixer design to state preparation,'' in \textit{Proc. IEEE Int. Conf. Quantum Comput. Eng. (QCE)}, 2020, pp. 72--82, doi: 10.1109/QCE49297.2020.00020.
\bibitem{Sherali1991} H. D. Sherali, T. B. Carter, and A. G. Hobeika, ``A location-allocation model and algorithm for evacuation planning under hurricane/flood conditions,'' \textit{Transp. Res. B, Methodol.}, vol. 25, no. 6, pp. 439--452, 1991, doi: 10.1016/0191-2615(91)90037-J.
\bibitem{Blekos2024} K. Blekos \textit{et al.}, ``A review on quantum approximate optimization algorithm and its variants,'' \textit{Phys. Rep.}, vol. 1068, pp. 1--66, 2024, doi: 10.1016/j.physrep.2024.03.002.
\bibitem{Fuchs2024} F. G. Fuchs and R. Pariente Bassa, ``LX-mixers for QAOA: Optimal mixers restricted to subspaces and the stabilizer formalism,'' \textit{Quantum}, vol. 8, Art. no. 1535, 2024, doi: 10.22331/q-2024-11-25-1535.
\bibitem{doCarmo2025} R. S. do Carmo, M. C. S. Santana, F. F. Fanchini, V. H. C. de Albuquerque, and J. P. Papa, ``Warm-starting QAOA with XY mixers: A novel approach for quantum-enhanced vehicle routing optimization,'' arXiv:2504.19934, 2025.
\bibitem{Gupta2026} A. Gupta and S. R. Hassan, ``Constraint-preserving QAOA for personnel rostering: Coverage-preserving and guarded-XY mixer constructions,'' arXiv:2607.09145, 2026.
\bibitem{Zhou2020} L. Zhou, S.-T. Wang, S. Choi, H. Pichler, and M. D. Lukin, ``Quantum approximate optimization algorithm: Performance, mechanism, and implementation on near-term devices,'' \textit{Phys. Rev. X}, vol. 10, no. 2, Art. no. 021067, 2020, doi: 10.1103/PhysRevX.10.021067.
\bibitem{Barkoutsos2020} P. K. Barkoutsos, G. Nannicini, A. Robert, I. Tavernelli, and S. Woerner, ``Improving variational quantum optimization using CVaR,'' \textit{Quantum}, vol. 4, Art. no. 256, 2020, doi: 10.22331/q-2020-04-20-256.
\bibitem{Egger2021} D. J. Egger, J. Mare\v{c}ek, and S. Woerner, ``Warm-starting quantum optimization,'' \textit{Quantum}, vol. 5, Art. no. 479, 2021, doi: 10.22331/q-2021-06-17-479.
\bibitem{Brandao2018} F. G. S. L. Brand\~{a}o, M. Broughton, E. Farhi, S. Gutmann, and H. Neven, ``For fixed control parameters the quantum approximate optimization algorithm's objective function value concentrates for typical instances,'' arXiv:1812.04170, 2018.
\bibitem{Galda2021} A. Galda, X. Liu, D. Lykov, Y. Alexeev, and I. Safro, ``Transferability of optimal QAOA parameters between random graphs,'' in \textit{Proc. IEEE Int. Conf. Quantum Comput. Eng. (QCE)}, 2021, pp. 171--180, doi: 10.1109/QCE52317.2021.00034.
\bibitem{Shaydulin2023} R. Shaydulin, P. C. Lotshaw, J. Larson, J. Ostrowski, and T. S. Humble, ``Parameter transfer for quantum approximate optimization of weighted MaxCut,'' \textit{ACM Trans. Quantum Comput.}, vol. 4, no. 3, Art. no. 19, 2023, doi: 10.1145/3584706.
\bibitem{Harrigan2021} M. P. Harrigan \textit{et al.}, ``Quantum approximate optimization of non-planar graph problems on a planar superconducting processor,'' \textit{Nat. Phys.}, vol. 17, pp. 332--336, 2021, doi: 10.1038/s41567-020-01105-y.
\bibitem{Bartschi2019} A. B\"{a}rtschi and S. Eidenbenz, ``Deterministic preparation of Dicke states,'' in \textit{Fundamentals of Computation Theory} (LNCS, vol. 11651). Cham, Switzerland: Springer, 2019, pp. 126--139, doi: 10.1007/978-3-030-25027-0\_9.
\bibitem{Vatan2004} F. Vatan and C. Williams, ``Optimal quantum circuits for general two-qubit gates,'' \textit{Phys. Rev. A}, vol. 69, no. 3, Art. no. 032315, 2004, doi: 10.1103/PhysRevA.69.032315.
\bibitem{Nielsen2010} M. A. Nielsen and I. L. Chuang, \textit{Quantum Computation and Quantum Information}, 10th anniversary ed. Cambridge, U.K.: Cambridge Univ. Press, 2010, doi: 10.1017/CBO9780511976667.
\bibitem{Powell1994} M. J. D. Powell, ``A direct search optimization method that models the objective and constraint functions by linear interpolation,'' in \textit{Advances in Optimization and Numerical Analysis}, S. Gomez and J.-P. Hennart, Eds. Dordrecht, The Netherlands: Kluwer, 1994, pp. 51--67, doi: 10.1007/978-94-015-8330-5\_4.
\bibitem{Huangfu2018} Q. Huangfu and J. A. J. Hall, ``Parallelizing the dual revised simplex method,'' \textit{Math. Program. Comput.}, vol. 10, no. 1, pp. 119--142, 2018, doi: 10.1007/s12532-017-0130-5.
\bibitem{Kirkpatrick1983} S. Kirkpatrick, C. D. Gelatt, and M. P. Vecchi, ``Optimization by simulated annealing,'' \textit{Science}, vol. 220, no. 4598, pp. 671--680, 1983, doi: 10.1126/science.220.4598.671.
\bibitem{Arute2019} F. Arute \textit{et al.}, ``Quantum supremacy using a programmable superconducting processor,'' \textit{Nature}, vol. 574, no. 7779, pp. 505--510, 2019, doi: 10.1038/s41586-019-1666-5.
\end{thebibliography}

\end{document}
